\begin{frame}{닫힌 형태의 의미: 대입하는 순간 끝}
  \begin{itemize}
    \item $\hat\mu_{y_i}$ = 샘플 $i$ 가 속한 \textbf{클래스의 평균}.
    \item 분모가 클래스별 $n_0,n_1$ 이 아니라 \textbf{전체 $N=n_0+n_1$}
      --- 바로 두 클래스의 편차를 \textbf{합쳐서} 평균내기.
    \item 1차원 예시의 ``편차 6개를 한꺼번에 나누던 것''과 같은 절차.
    \item 로지스틱회귀: 경사하강으로 $w,b$ 를 \textbf{반복 탐색}
      vs GDA: 수식에 \textbf{대입하는 순간 파라미터 끝}.
    \item Week 2.2 정규방정식의 ``한 번에 풀기''와 같은 계열.
    \item {\footnotesize (세세한 차이: 통계학에선 $N{-}d{-}1$ 로 나누어
      편향을 없애는 버전도 있지만, 학습엔 위 MLE 버전이 표준,
      데이터가 크면 차이 없음.)}
  \end{itemize}
\end{frame}
